\documentclass[10pt,a4paper]{amsart}
\usepackage[margin=19mm]{geometry}
\usepackage{amsmath,amssymb,mathtools}
\usepackage[scr=rsfs,cal=euler]{mathalfa}
\usepackage{xfrac}
\usepackage[unicode,hidelinks,bookmarks=false]{hyperref}
\newcommand{\ie}{i.e.\ }
\newcommand{\cf}{cf.\ }
\newcommand{\resp}{respectively\ }
\newcommand{\Subsection}{\subsection}
\providecommand{\nobreakdash}{\nobreak\hbox{-}}
\setlength{\emergencystretch}{2em}
\multlinegap0pt
\usepackage{xcolor}
\usepackage{algorithm}\usepackage{algpseudocode}
\usepackage{amssymb}
\usepackage{mathtools}
\newtheorem{lemma}{Lemma}
\title{Regret Bounds for Expected Improvement Algorithms in Gaussian Process Bandit Optimization}
\author{Hung Tran-The, Sunil Gupta, Santu Rana, Svetha Venkatesh}
\date{Source: arXiv:2203.07875v1 (CC BY 4.0)}
\begin{document}
\maketitle

\noindent\emph{Source and licence.} Regret Bounds for Expected Improvement Algorithms in Gaussian Process Bandit Optimization. Version arXiv:2203.07875v1 (CC BY 4.0), distributed by arXiv under the Creative Commons Attribution 4.0 International licence, http://creativecommons.org/licenses/by/4.0/, as recorded in the arXiv licence metadata for this exact version and shown on the abstract page https://arxiv.org/abs/2203.07875v1. Full licence notice: \texttt{licenses/arxiv-2203.07875-CC-BY-4.0.txt}.\par\smallskip

\noindent\textit{Editorial scope.} Counted unit: the article's lemma lem3.5 (source0.tex lines 787--793). The article's complete original proof of it is delivered with it (source0.tex lines 794--805, one \texttt{proof} environment of 12 lines). No mathematics is written, repaired or re-derived: the statement and the proof are the article's own lines, in the article's own notation, and no retained line is substituted or re-flowed.  No further source-local context is needed: every cross-reference of every retained line resolves inside the two delivered spans.  Nothing was omitted from any retained span, so the package carries no omission record.  No retained line contains a citation, so the wrapper carries no bibliography and no bibliography entry.  No retained line contains a figure environment, an image-inclusion command or drawing code, so no image is delivered and the package declares no asset dependency.  Editorial formatting only, in addition: an AMS \texttt{amsart} preamble that restates the article's own declarations and macros used by the retained lines, namely the environments \texttt{lemma} on separate counters, together with the standard packages those lines need.  The retained environments print in this document as Lemma 1 in order of appearance, whereas the article prints the same items with the numbering of its own full document.  The complete native source of the article as distributed in the arXiv e-print for this exact version is bundled unchanged as \texttt{sources/123264/source0.tex}.
\begin{lemma}
Let $\omega_T = \sqrt{\text{ln}(T)\text{ln} \text{ln}(T)}$. Given a set $A \in \mathcal A_t$ and assume that there exists a $1 \le \phi \le T$ such that $A \in \mathcal A_{\tau}$. Let $\phi'(A) = \text{max}\{1 \le t \le T: A \in \mathcal A_t\}$. Then with probability at least $1 - \delta$, we have
$$\sum_{i=1}^{n(A)}r^A_i \le \tilde{\mathcal O}(\sqrt{n(A)\gamma^A_{\phi'(A)}}),$$
where $n(A) = |\mathcal D_T \cap A|$ which is defined as above. The notation $\tilde{\mathcal O}$ is a
variant of $\mathcal O$, where log factors are suppressed.
\label{lem3.5}
\end{lemma}

\begin{proof}
Following the Improved-GP-EI algorithm, each $A \in \mathcal A_t$ is fitted by an independent Gaussian process with $N(A)$ observations and the sampled points in $\mathcal D_t^A$: $(x_1^A,y_1^A), ..., (x_{N(A)}^A,y_{N(A)}^A)$. Therefore, we can apply the results in Section B to the set $A$.

Similar to the proofs of Lemma 6, Lemma 7, and Lemma 8, we obtain an upper bound on $r^A_i = f(x^*) - f(x^{A+}_i)$ as follows.
$$r^A_i \le (\frac{\tau(\frac{\beta^{A}_{t_i}}{\omega_T})}{ \tau(-\frac{\beta^{A}_{t_i}}{\omega_T})} + \sqrt{2\pi}) \text{max}\{0, f(x^A_{i+1}) - \mu^{A+}_{t_i}\} + (\frac{\tau(\frac{\beta^{A}_{t_i}}{\omega_T})}{ \tau(-\frac{\beta^{A}_{t_i}}{\omega_T})}(\beta^{A}_{t_i} + \omega_T) + \sqrt{2\pi}(3\beta^{A}_{t_i} + \omega_T))\sigma^A_{t_i}(x^A_{i+1})),$$
where $\mu^{A+}_i = \text{max}_{x^A_j \in \mathcal D_T \cap A, 1\le j \le i} \mu_{t_i}(x^A_j)$. Recall that $x^{A+}_i = \text{argmax}_{x^A_j \in \mathcal D_T \cap A, 1\le j \le i} f(x^A_j)$.

Using Lemma 21, we can bound the ratio $\frac{\tau(\frac{\beta^A_{t_i}}{\omega_T})}{ \tau(-\frac{\beta^A_{t_i}}{\omega_T})}$ by a constant which is independent of $1 \le t\le T$ and $T$. Thus, by the proofs similar to those of Lemma 14 and Lemma 15, we obtain an upper bound on the sum $\sum_{i=1}^{n(A)}r^A_i$ as
$$\sum_{i=1}^{n(A)}r^A_i \le \mathcal O(\beta^A_{t_{N(A)}}\sqrt{n(A)\gamma^A_{\phi'(A)}}).$$
Using the proof similar to Lemma 21 as above, there exists constant $C_1$ such that $\beta^A_{t_{N(A)}} \le C_1\omega_T = \Theta(\sqrt{\text{ln}(T)\text{ln} \text{ln}(T)})$. We can remove this log component from the upper bound. Finally, we have that
$$\sum_{i=1}^{n(A)}r^A_i \le \tilde{\mathcal O}(\sqrt{n(A)\gamma^A_{\phi'(A)}}).$$
\end{proof}

\end{document}
